% Part: lambda-calculus
% Chapter: syntax
% Section: abbreviated-syntax

\documentclass[../../../include/open-logic-section]{subfiles}

\begin{document}

\olfileid{lam}{syn}{abb}
\olsection{النحو المختصر}

قد تطول كتابة الحدود على مقتضى~\olref[trm]{defn:term} وتثقل؛
فلذلك نستعمل نحوًا أوجز، بشرط أن تبقى قراءة كل حد وحيدة.
ومن الاصطلاحات الشائعة في ذلك ما يسمى \emph{الحدود المختصرة}، وقواعده الآتية.

\begin{enumerate}
\item إذا أسقطت الأقواس، جمع التطبيق من جهة اليسار إلى اليمين.
  فإذا كانت~$M$ و$N$ و$P$ و~$Q$ حدودًا، فالكتابة~$MNPQ$ تختصر~$(((MN)P)Q)$.
\item وإذا أسقطت الأقواس، امتد مجال تجريد لامبدا إلى أقصى ما تسمح به الكتابة.
  ولذلك تقرأ~$\lambd[x][MNP]$ بمعنى~$(\lambd[x][MNP])$.
\item ويجوز أن يجمع رمز لامبدا تجريد عدة متغيرات؛ فمثلًا
  $\lambd[xyz][M]$ كتابة مختصرة للحد
  $\lambd[x][\lambd[y][\lambd[z][M]]]$.
\end{enumerate}

ومن أمثلة هذه القواعد أن
\[
\lambd[xy][xxyx \lambd[z][xz]]
\]
تقوم مقام
\[
(\lambd[x][(\lambd[y][((((xx)y)x)(\lambd[z][(xz)]))])]).
\]

\begin{prob}
اكتب الحد المختصر~$\lambd[g][(\lambd[x][g (x x)])
  \lambd[x][g (x x)]]$ بإثبات جميع الأقواس.
\end{prob}

\end{document}

